% =====================================================================
\chapter{Related Work}
\label{ch:relatedwork}
% =====================================================================

\section{Firmware Rehosting and Emulation}
\label{sec:rw:rehosting}

The systematisation by Fasano et al.~\cite{fasano2021sok} and the
survey by Wright et al.~\cite{wright2021challenges} map the design
space of rehosting and identify \emph{fidelity} as its central
challenge. Both works focus predominantly on \emph{execution} fidelity
(does the firmware run at all, and does it take the same paths?) and
\emph{memory/peripheral} fidelity; temporal fidelity is acknowledged
but not quantified. This thesis fills that gap with hardware-grounded
measurements.

Firmadyne~\cite{chen2016firmadyne} pioneered large-scale emulation of
Linux-based firmware; HALucinator~\cite{clements2020halucinator}
replaces hardware-abstraction-layer functions with host-side handlers;
P$^2$IM~\cite{feng2020p2im}, Pretender~\cite{gustafson2019pretender},
and $\mu$Emu~\cite{zhou2021uemu} infer peripheral behaviour
automatically to keep firmware running. All of these systems evaluate
success as ``firmware executes correctly'', with no claim about
\emph{when} events happen. Muench et al.~\cite{muench2018you}
demonstrate that even memory-corruption effects differ between
emulated and physical embedded targets---an early warning that
emulation platforms are not behaviourally transparent. Our work extends
this observation to the time domain.

\section{Timing-Accurate Simulation}
\label{sec:rw:timing}

At the opposite end of the fidelity spectrum, cycle-accurate simulators
model microarchitectural state per clock cycle. gem5~\cite{binkert2011gem5}
provides cycle-approximate CPU models for architecture research, and
commercial virtual platforms such as Simics~\cite{magnusson2002simics}
offer configurable timing models. SystemC/TLM-2.0~\cite{systemc} defines
loosely- and approximately-timed modelling styles in which
transaction-level models annotate delays. These approaches achieve high
temporal fidelity at the cost of simulation speed (typically
$10^2$--$10^5\times$ slower than real time) and modelling effort; they
are rarely practical for full-firmware CI workflows. The Temporal
Fidelity Layer targets the middle ground: statistical (rather than
cycle-exact) timing realism at near-native emulation speed.

\section{Emulator Time Management}
\label{sec:rw:quantum}

QEMU's icount mode~\cite{bellard2005qemu} maps executed instructions to
virtual time at a fixed rate, analogous to Renode's MIPS-based model,
and shares its limitation: instruction counts, not microarchitectural
cycles, drive the clock. Renode's quantum-based synchronisation of
multiple machines~\cite{renode} enables deterministic multi-node
simulation but aligns all cross-machine communication to quantum
boundaries. The evaluation in \cref{sec:eval:quantum} quantifies the
fidelity--performance trade-off of this quantum for the first time on
a hardware-validated workload.

\section{Positioning of this Thesis}
\label{sec:rw:positioning}

To the best of our knowledge, this thesis provides the first
systematic, statistically evaluated comparison of RTOS-level timing
primitives between a physical Cortex-M platform and its rehosted
counterpart in Renode, together with a non-invasive mitigation layer.
Unlike cycle-accurate simulation, TFL does not attempt to reproduce
\emph{individual} event times; it aims to reproduce hardware-like
timing \emph{distributions}---sufficient for testing timing-robustness
of firmware---while preserving emulation speed and determinism controls.
